\documentclass[12pt]{article}

\usepackage[utf8]{inputenc}
\usepackage[T1]{fontenc}
\usepackage{lmodern}
\usepackage[a4paper,margin=1in]{geometry}
\usepackage{setspace}
\usepackage{booktabs}
\usepackage{amsmath}
\usepackage{graphicx}
\usepackage{hyperref}
\usepackage{enumitem}
\usepackage{microtype}
\usepackage{xurl}

\setstretch{1.15}

\title{Object Relations and Identity Features in Women with Bulimia Nervosa:\\An Exploratory Clinical Pilot Study}
\author{Evgeniya S. Profatilova\thanks{Correspondence: \href{mailto:e.profatilova@gmail.com}{e.profatilova@gmail.com}}\\
Independent Researcher, Limassol, Cyprus}
\date{March 2026}

\begin{document}

\maketitle

\begin{quote}
\small
\noindent\textit{SSRN preprint (working paper version).} This manuscript is an English-language adaptation of a 2025 master's thesis dataset and is shared for scholarly discussion. The underlying thesis was completed at Lomonosov Moscow State University.
\normalsize
\end{quote}

\begin{abstract}
\noindent
\textbf{Background:} Psychodynamic models of bulimia nervosa emphasize disturbances in object relations and identity organization, but quantitative evidence from clinical samples remains limited. \\
\textbf{Objective:} To examine associations between object-relations characteristics and identity indicators in women with bulimia nervosa. \\
\textbf{Methods:} This exploratory cross-sectional pilot study included 28 women (16--40 years) with clinically confirmed bulimia nervosa (ICD-10 F50.2), recruited from a specialized eating-disorders center. Participants completed the Bell Object Relations Inventory (BORI) and Erwin Identity Scale (EIS-III). Analyses included descriptive statistics, Pearson correlations, and multiple linear regressions. \\
\textbf{Results:} Insecure attachment ($M=6.57$, $SD=3.08$) and alienation ($M=6.07$, $SD=3.64$) were the most pronounced object-relations features. Identity domains were strongly interrelated (confidence with body/appearance: $r=0.91$, $p<0.001$; confidence with sexual-identity indicators: $r=0.87$, $p<0.001$). Cross-domain analyses showed that higher object-relations disturbance was associated with less favorable identity indicators (e.g., alienation with confidence: $r=-0.75$, $p<0.01$). Regression models explained 58\% of variance in confidence, 62\% in sexual-identity indicators, and 39\% in body/appearance indicators. \\
\textbf{Conclusions:} In this pilot clinical sample, object-relations disturbances were robustly associated with identity-related difficulties. Findings are consistent with integrative psychodynamic case formulation but should be interpreted as exploratory given design and sample-size constraints.
\end{abstract}

\noindent\textbf{Keywords:} bulimia nervosa; object relations; identity; insecure attachment; alienation; psychodynamic psychology.

\section{Introduction}
Bulimia nervosa is associated with recurrent binge-eating episodes, compensatory behaviors, emotional dysregulation, and substantial psychosocial burden (Bruch, 1973; Smyth et al., 2007). Beyond symptom counts, current clinical literature highlights chronic shame, interpersonal insecurity, and dependence on external evaluation as key maintaining processes (Zerbe, 2016).

Psychodynamic approaches conceptualize these difficulties through two linked domains: object relations and identity organization. In this framework, problematic internalized relationship patterns (e.g., alienation, insecure attachment) may be associated with low confidence, unstable self-representation, and body-related self-evaluation. Prior empirical work has reported links between eating-disorder pathology and both attachment disturbances and identity impairment (Becker et al., 1987; Sacchetti et al., 2019), but integrated analyses in single clinical samples remain relatively uncommon.

The present study addresses that gap using standardized instruments in a clinical sample of women with bulimia nervosa. The analysis focused on three questions: (1) Which object-relations features are most pronounced? (2) How are identity domains interrelated? (3) How strongly are object-relations features associated with identity indicators?

The primary hypothesis was that higher levels of object-relations disturbance (especially alienation and insecure attachment) would be associated with less favorable identity indicators.

\section{Methods}

\subsection{Design}
This was an exploratory cross-sectional pilot study based on the original master's-thesis dataset. Reporting was structured to align with core transparency principles used in observational research.

\subsection{Setting and participants}
Participants were recruited from a specialized eating-disorders treatment center in Moscow. The sample included 28 women receiving treatment for clinically confirmed bulimia nervosa (ICD-10: F50.2), aged 16--40 years. No control group was included.

\subsection{Measures}
\textbf{Bell Object Relations Inventory (BORI).} Four domains were analyzed: Alienation, Insecure Attachment, Egocentricity, and Social Incompetence.

\textbf{Erwin Identity Scale (EIS-III).} Three domains were analyzed: Confidence, Sexual-Identity indicators, and Body/Appearance indicators.

\subsection{Procedure}
Participants completed questionnaires in a confidential clinical setting. The source document reports informed consent prior to participation.

\subsection{Statistical analysis}
Descriptive statistics were calculated for all domains. Pearson correlations were used to evaluate associations within BORI, within EIS-III, and across instruments. Multiple linear regression models examined the contribution of BORI domains to each EIS-III outcome. All inferential tests were two-tailed with a significance threshold of $p<0.05$. Because of the pilot design and sample size, all inferential findings are interpreted as exploratory rather than confirmatory.

\subsection{Transparency and reporting statements}
No formal preregistration record was reported in the source thesis. No a priori power analysis was reported. The current manuscript reports the available model outputs from the source dataset and explicitly labels inferential findings as exploratory.

\section{Results}

\subsection{Descriptive statistics}
Table~\ref{tab:bori_desc} presents BORI descriptive values. Insecure attachment and alienation were the highest-scoring domains.

\begin{table}[htbp]
\centering
\caption{Descriptive statistics for BORI domains}
\label{tab:bori_desc}
\begin{tabular}{lcccc}
\toprule
Domain & Mean & SD & Min & Max \\
\midrule
Alienation & 6.07 & 3.64 & 0 & 11 \\
Insecure attachment & 6.57 & 3.08 & 1 & 12 \\
Egocentricity & 3.86 & 1.35 & 1 & 6 \\
Social incompetence & 4.82 & 2.28 & 0 & 8 \\
\bottomrule
\end{tabular}
\end{table}

Table~\ref{tab:eis_desc} presents EIS-III descriptive values and deviations from the scale-specific neutral points used in the source analysis.

\begin{table}[htbp]
\centering
\caption{Descriptive statistics for EIS-III domains}
\label{tab:eis_desc}
\begin{tabular}{lccc}
\toprule
Domain & Mean & SD & Deviation from neutral point \\
\midrule
Confidence & 68.25 & 14.91 & $-3.75$ \\
Sexual-identity indicators & 53.46 & 10.52 & $-0.54$ \\
Body/appearance indicators & 54.46 & 11.51 & $+3.46$ \\
\bottomrule
\end{tabular}
\end{table}

\subsection{Correlation findings}
Within BORI domains, the strongest associations were alienation with social incompetence ($r=0.74$, $p<0.001$) and insecure attachment with social incompetence ($r=0.68$, $p<0.001$). Alienation also correlated with insecure attachment ($r=0.52$, $p=0.0048$).

Within EIS-III domains, all pairwise associations were strong and significant: confidence with sexual-identity indicators ($r=0.87$, $p<0.001$), confidence with body/appearance indicators ($r=0.91$, $p<0.001$), and sexual-identity indicators with body/appearance indicators ($r=0.88$, $p<0.001$).

Cross-domain results are summarized in Table~\ref{tab:crosscorr}. The strongest association was between alienation and confidence ($r=-0.75$, $p<0.01$).

\begin{table}[htbp]
\centering
\caption{Selected cross-domain correlations between BORI and EIS-III}
\label{tab:crosscorr}
\small
\begin{tabular}{lccc}
\toprule
BORI domain & Confidence & Sexual identity & Body and appearance \\
\midrule
Alienation & $-0.75$ ($p<0.01$) & $-0.68$ ($p<0.01$) & $-0.62$ ($p<0.01$) \\
Insecure attachment & $-0.72$ ($p<0.01$) & $-0.61$ ($p<0.01$) & $-0.56$ ($p<0.01$) \\
Egocentricity & $-0.64$ ($p<0.01$) & $-0.57$ ($p<0.01$) & $-0.50$ ($p<0.05$) \\
Social incompetence & $-0.66$ ($p<0.01$) & $-0.59$ ($p<0.01$) & $-0.53$ ($p<0.05$) \\
\bottomrule
\end{tabular}
\end{table}

\subsection{Regression findings}
Model-level fit was as follows: confidence $R^2=0.58$, $F(4,23)=7.94$, $p<0.001$; sexual-identity indicators $R^2=0.62$, $F(4,23)=9.38$, $p<0.001$; body/appearance indicators $R^2=0.391$, $F(4,23)=3.69$, $p=0.018$.

Thus, all three models were statistically significant at the conventional $\alpha=0.05$ level, while the body/appearance model showed a weaker level of evidence compared with the other two outcomes.

Table~\ref{tab:regression} summarizes standardized coefficients ($\beta$) from source models.

\begin{table}[htbp]
\centering
\caption{Regression coefficients for EIS-III outcomes predicted by BORI domains}
\label{tab:regression}
\small
\resizebox{\textwidth}{!}{%
\begin{tabular}{lccc}
\toprule
Predictor & Confidence ($\beta$, $p$) & Sexual identity ($\beta$, $p$) & Body and appearance ($\beta$, $p$) \\
\midrule
Alienation & $-0.45$, 0.002 & $0.38$, 0.004 & $0.412$, $<0.001$ \\
Insecure attachment & $-0.38$, 0.002 & $0.45$, 0.001 & $0.265$, 0.024 \\
Egocentricity & $-0.11$, 0.366 & $-0.22$, 0.058 & $0.128$, 0.187 \\
Social incompetence & $-0.27$, 0.012 & $0.31$, 0.002 & $0.352$, 0.003 \\
\bottomrule
\end{tabular}
}
\end{table}

A directionality inconsistency should be noted: for two EIS-III outcomes, regression-coefficient signs in source tables differ from the sign pattern in bivariate correlations. Because item-level recoding and full scoring-direction documentation were unavailable in the source material, directional interpretation of regression coefficients should be treated cautiously.

\section{Discussion}
This exploratory pilot analysis indicates that women with bulimia nervosa in this clinical sample showed pronounced insecure attachment and alienation, and that these object-relations features were strongly associated with identity-related indicators.

The findings are compatible with psychodynamic formulations in which interpersonal insecurity, social withdrawal, and fragile self-representation mutually reinforce symptom maintenance (Zerbe, 2016). In particular, the magnitude and consistency of cross-domain correlations suggest that interpersonal/object-relations disturbances and identity processes are tightly linked at the clinical-phenomenological level, which is also consistent with prior case-control observations on related interpersonal-mentalizing deficits (Sacchetti et al., 2019).

At the same time, interpretation should remain conservative. The study was not designed to establish causality, and coefficients should not be interpreted as directional mechanisms.

\subsection{Clinical implications (exploratory)}
The results support further development and testing of integrative treatment targets that address:
\begin{itemize}[leftmargin=*]
    \item alienation and interpersonal avoidance,
    \item attachment-related anxiety and dependence on external validation,
    \item confidence and self-representation stability,
    \item body-related self-evaluation in the context of broader identity work.
\end{itemize}

\subsection{Limitations}
The study has substantial limitations relevant for inference strength:
\begin{itemize}[leftmargin=*]
    \item small sample size ($n=28$),
    \item single-site clinical recruitment,
    \item absence of a control group,
    \item cross-sectional design,
    \item reliance on self-report instruments,
    \item no formal symptom-severity covariates in the main models,
    \item incomplete documentation of scoring direction details for secondary interpretation,
    \item no ethics-committee identifier reported in available source text.
\end{itemize}

\subsection{Future research}
Future work should use larger multi-site samples, include control/comparison groups, pre-register analytic plans, and apply longitudinal designs with symptom-severity and treatment-process indicators.

\section{Conclusion}
In this exploratory clinical pilot sample, object-relations disturbances were strongly associated with identity-related difficulties in women with bulimia nervosa. The pattern is clinically meaningful and theoretically coherent, but it requires confirmation in larger, controlled, and longitudinal studies.

\section*{Declarations}
\textbf{Ethics statement.} The source thesis reports informed consent. An ethics-committee approval number was not available in the analyzed document.

\textbf{Preregistration.} No preregistration identifier was reported in the source document.

\textbf{Funding.} No external funding was reported in the source document.

\textbf{Conflicts of interest.} No conflicts of interest were reported in the source document.

\textbf{Data availability.} De-identified participant-level data were not publicly available in the source document. The computational verification script is publicly available at Zenodo: \href{https://doi.org/10.5281/zenodo.18986931}{https://doi.org/10.5281/zenodo.18986931} (version v1.0.0).

\section*{References}
\begin{enumerate}[leftmargin=*]
    \item Becker, A. E., Grinspoon, S. K., Klibanski, A., \& Herzog, D. B. (1987). Eating disorders and object relations in women. \textit{Journal of Clinical Psychology, 43}(1), 92--95. doi:10.1002/1097-4679(198701)43:1<92::AID-JCLP2270430113>3.0.CO;2-Z
    \item Bruch, H. (1973). \textit{Eating Disorders: Obesity, Anorexia Nervosa, and the Person Within}. New York: Basic Books.
    \item Sacchetti, S., Robinson, P., Bogaardt, A., Clare, A., Ouellet-Courtois, C., Luyten, P., Bateman, A., \& Fonagy, P. (2019). Reduced mentalizing in patients with bulimia nervosa and features of borderline personality disorder: A case-control study. \textit{BMC Psychiatry, 19}, Article 134. doi:10.1186/s12888-019-2112-9
    \item Smyth, J. M., Wonderlich, S. A., Heron, K. E., Sliwinski, M. J., Crosby, R. D., Mitchell, J. E., \& Engel, S. G. (2007). Daily and momentary mood and stress are associated with binge eating and vomiting in bulimia nervosa patients in the natural environment. \textit{Journal of Consulting and Clinical Psychology, 75}(4), 629--638. doi:10.1037/0022-006X.75.4.629
    \item Zerbe, K. (2016). Psychodynamic issues in the treatment of binge eating: Working with shame, secrets, no-entry, and false body defenses. \textit{Clinical Social Work Journal, 44}(1), 8--17. doi:10.1007/s10615-015-0559-9
\end{enumerate}

\end{document}